% Slides de defesa — tese curta (PT-PT).
% ~15 minutos. Regra: um slide, uma ideia. O que está no slide é o que se aponta; o resto diz-se.
% O penúltimo frame é a DEMONSTRAÇÃO (gravação da app).

\documentclass[aspectratio=169]{beamer}
\usetheme{Madrid}
\usecolortheme{seahorse}
\usepackage[T1]{fontenc}
\usepackage[portuguese,shorthands=off]{babel}
\usepackage{booktabs}
\usepackage{tikz}
\usetikzlibrary{positioning, arrows.meta, calc, decorations.pathreplacing}
\tikzset{every node/.append style={execute at begin node={\hyphenpenalty=10000\exhyphenpenalty=10000\relax}}}
\graphicspath{{../../thesis/figures/}}
\setbeamertemplate{navigation symbols}{}
\setbeamertemplate{caption}{\raggedright\insertcaption\par}

\title[InvestiGator]{Alertas Financeiros Explicáveis para Investidores de Retalho}
\author[H.\ Santos]{Henrique José da Silva Santos \\
  \footnotesize Orientador: Prof.\ Luís Gomes \quad Coorientador: Rafael Silva}
\institute[ISEP]{MEIA, Instituto Superior de Engenharia do Porto}
\date{\footnotesize Defesa}

\begin{document}

\begin{frame}\titlepage\end{frame}

% ══ 1. O PROBLEMA ═══════════════════════════════════════════════════════════
\begin{frame}{Uma pessoa abre a aplicação e vê isto}
\begin{center}
\begin{tikzpicture}
  \node[draw, rounded corners=6pt, minimum width=5cm, minimum height=2.6cm, fill=black!4] (t) {};
  \node[font=\Large\bfseries] at ([yshift=6mm]t.center) {NVDA};
  \node[font=\Huge\bfseries, text=red!65!black] at ([yshift=-6mm]t.center) {$-4{,}1\%$};
\end{tikzpicture}
\end{center}
\vspace{2mm}
\begin{center}
\large E fica com \textbf{três perguntas}, sempre estas e sempre por esta ordem.
\end{center}
\end{frame}

\begin{frame}{As três perguntas}
\begin{enumerate}
  \item[\textbf{1.}] \textbf{Isto é invulgar?}\\
        \footnotesize 4\% numa empresa calma é muito. Noutra, é terça-feira.\\[3mm]
  \item[\textbf{2.}] \textbf{Foi a empresa, ou foi o mercado?}\\
        \footnotesize Se caiu tudo, a história é outra.\\[3mm]
  \item[\textbf{3.}] \textbf{Já aconteceu antes, e o que se seguiu?}\\
        \footnotesize Houve dias parecidos? O que fez o preço depois?
\end{enumerate}
\vspace{3mm}
\begin{center}
\textbf{As aplicações gratuitas não respondem a nenhuma delas.}
\end{center}
\end{frame}

\begin{frame}{O que existe hoje}
\centering
\small
\begin{tabular}{lcccl}
\toprule
& \textbf{invulgar?} & \textbf{empresa ou mercado?} & \textbf{já aconteceu?} & \textbf{custo} \\
\midrule
App da corretora & --- & --- & --- & incluído \\
Agregador de notícias & --- & --- & --- & grátis \\
\emph{Robo-advisor} & --- & --- & --- & \% da carteira \\
Terminal profissional & \checkmark & \checkmark & \checkmark & milhares/ano \\
\midrule
\textbf{Este trabalho} & \checkmark & \checkmark & \checkmark & \textbf{grátis} \\
\bottomrule
\end{tabular}
\vspace{5mm}

\begin{center}
A linha que responde às três \textbf{já existe}. É a mais cara de todas.\\
\textbf{O problema não é científico: é de acesso.}
\end{center}
\end{frame}

% ══ 2. O SISTEMA ════════════════════════════════════════════════════════════
\begin{frame}{O que o sistema faz}
\begin{center}
\begin{tikzpicture}[
  font=\small, >={Stealth[length=2.2mm]},
  gat/.style={draw, rounded corners, fill=black!8, align=center, text width=2.9cm, minimum height=0.9cm},
  sis/.style={draw, rounded corners, fill=black!3, align=center, text width=2.4cm, minimum height=1.5cm},
  sa/.style={draw, rounded corners, fill=black!3, align=center, text width=3.4cm, minimum height=1.5cm},
]
  \node[gat] (g1) {movimento invulgar};
  \node[gat, below=6mm of g1] (g2) {notícia relevante};
  \node[sis] (s) at ($(g1)!0.5!(g2) + (3.9,0)$) {\textbf{InvestiGator}};
  \node[sa, right=1.1cm of s] (o) {alerta explicado\\[2pt]
    \scriptsize o quê · porquê\\ \scriptsize fontes · casos passados};
  \draw[->] (g1) -- (s); \draw[->] (g2) -- (s); \draw[->] (s) -- (o);
\end{tikzpicture}
\end{center}
\vspace{4mm}
\begin{center}
\large E há uma coisa que ele \textbf{recusa} fazer.
\end{center}
\end{frame}

\begin{frame}{Não prevê preços}
\begin{center}
\Large
Nunca diz que vai subir.\\[2mm]
Nunca diz para comprar.\\[2mm]
Nunca mostra um alvo de preço.
\end{center}
\vspace{6mm}
\begin{center}
\normalsize
Parece uma limitação. É o contrário:\\
\textbf{é o que torna tudo o resto verificável.}\\[2mm]
\footnotesize Uma afirmação sobre o que já aconteceu confere-se contra o registo.\\
Uma previsão não se confere.
\end{center}
\end{frame}

% ══ 3. AS TRÊS TÉCNICAS ═════════════════════════════════════════════════════
\begin{frame}{Técnica 1: é invulgar?}
Comparar o dia de hoje com \textbf{os 20 dias anteriores da mesma ação}.
\begin{center}
\begin{tikzpicture}[font=\footnotesize]
  \foreach \i in {0,...,19} { \fill[black!18] (\i*0.28,0) rectangle (\i*0.28+0.2,0.45); }
  \fill[orange!70] (20*0.28,0) rectangle (20*0.28+0.2,1.0);
  \draw[decorate, decoration={brace, amplitude=4pt}, black!55] (0,-0.12) -- (19*0.28+0.2,-0.12)
    node[midway, below=5pt, black!65] {os 20 dias anteriores};
  \draw[decorate, decoration={brace, amplitude=4pt}, orange!75!black]
    (20*0.28,-0.12) -- (20*0.28+0.2,-0.12) node[midway, below=5pt, text=orange!60!black] {hoje};
\end{tikzpicture}
\end{center}
\vspace{1mm}
\begin{center}
$z = \dfrac{\text{movimento de hoje} - \text{média}}{\text{desvio-padrão}}$
\end{center}
\vspace{2mm}
\footnotesize
$z=+7$ quer dizer o mesmo numa empresa calma e numa volátil.\\
\textbf{O dia julgado não entra no cálculo da sua própria norma.} Sem isto, seria falsear.
\end{frame}

\begin{frame}{Técnica 2: já aconteceu antes?}
Procurar por palavras falha nos \textbf{dois} sentidos:
\vspace{2mm}
\begin{center}
\begin{tikzpicture}[font=\footnotesize]
  \fill[black!70] (0.6,1.5) circle (2pt) node[right=2pt] {``bate estimativas''};
  \fill[black!70] (1.1,2.0) circle (2pt) node[right=2pt] {``supera previsões''};
  \draw[black!35, dashed] (0.85,1.75) circle (0.75);
  \node[black!50, font=\scriptsize\itshape] at (0.9,2.75) {mesmo sentido, palavras diferentes};
  \fill[orange!75!black] (6.6,0.8) circle (2pt) node[right=2pt] {``bate recorde de despedimentos''};
  \draw[orange!55, dashed] (6.6,0.8) circle (0.55);
  \node[orange!60!black, font=\scriptsize\itshape] at (6.9,1.7) {palavra igual, sentido diferente};
\end{tikzpicture}
\end{center}
\vspace{1mm}
Cada manchete passa a \textbf{384 números}: uma coordenada de significado.\\
Compara-se pelo \textbf{ângulo} entre elas.
\vspace{2mm}

Depois, para cada caso passado, mede-se o que o preço fez a $+1$, $+3$ e $+5$ dias.\\
\footnotesize \textbf{É uma medição do passado, não uma previsão.}
\end{frame}

\begin{frame}{Técnica 3: merece avisar?}
Aqui não há fórmula óbvia. É o sítio para \textbf{aprender dos dados}.
\vspace{3mm}
\begin{itemize}
  \item \textbf{79\,753} exemplos que construí a partir de notícias e preços (2018--2023)
  \item Pergunta: \emph{depois desta notícia houve um movimento anormalmente grande?}
  \item \textbf{Em valor absoluto}, nunca em que direção
  \item Divisão \textbf{cronológica} com embargo; probabilidades calibradas
\end{itemize}
\vspace{4mm}
\begin{center}
\footnotesize O modelo, o seu treino e as suas métricas ficam guardados como artefactos.\\
\textbf{É o meu modelo, treinado nos meus dados.}
\end{center}
\end{frame}

% ══ 4. UM ALERTA REAL ═══════════════════════════════════════════════════════
\begin{frame}{Um alerta real, do princípio ao fim}
\footnotesize Meta, 12 de julho de 2026. \emph{``Mark Zuckerberg Said Meta's AI Bets `Haven't Come
to Fruition Yet' as Shares Fell 5\%''}
\vspace{2mm}
\centering
\scriptsize
\begin{tabular}{rll}
\toprule
& \textbf{etapa} & \textbf{o que aconteceu} \\
\midrule
2 & nomeia a empresa? & ``Meta'', ``Zuckerberg'' \checkmark \\
3 & é recente? & menos de 2 dias \checkmark \\
4 & casos parecidos & MSFT $0{,}46$ · NVDA $0{,}51$ · NVDA $0{,}46$ \\
5 & evidência forte? & melhor $0{,}51 \ge 0{,}45$ \checkmark \\
6 & merece avisar? & $54\% \ge 50\%$ \checkmark \\
9 & enviado & ao canal, e guardado \\
\bottomrule
\end{tabular}
\vspace{3mm}

\begin{center}
\footnotesize Os desfechos foram \textbf{em sentidos opostos}: $+2{,}70\%$ · $-3{,}87\%$ · $+5{,}05\%$\\
Por isso mostram-se \textbf{um a um}: uma média esconderia isso.
\end{center}
\end{frame}

% ══ 5. RESULTADOS ═══════════════════════════════════════════════════════════
\begin{frame}{Resultado 1: deteção \textcolor{green!50!black}{\textbf{sim}}}
Um bom detetor deve disparar a um ritmo \textbf{parecido} em empresas diferentes.
\vspace{2mm}
\centering
\begin{tabular}{lccc}
\toprule
& \textbf{mín} & \textbf{máx} & \textbf{amplitude} \\
\midrule
\emph{z}-score & $0{,}016$ & $0{,}031$ & $\mathbf{0{,}015}$ \\
Limiar fixo de 3\% & $0{,}009$ & $0{,}353$ & $0{,}344$ \\
\bottomrule
\end{tabular}
\vspace{4mm}

\begin{center}
Mais de \textbf{20 vezes} mais consistente.\\[2mm]
\small E ganha também a dois métodos aprendidos:
$F_1$ $\mathbf{0{,}530}$ contra $0{,}269$ e $0{,}280$.
\end{center}
\end{frame}

\begin{frame}{Resultado 2: precedentes \textcolor{green!50!black}{\textbf{sim}}}
\centering
\begin{tabular}{lc}
\toprule
\textbf{método} & \textbf{precisão@5} \\
\midrule
\textbf{Por significado (o método)} & $\mathbf{0{,}514}$ \\
Por palavras em comum & $0{,}346$ \\
Ao acaso (o chão) & $0{,}240$ \\
As mais recentes & $0{,}126$ \\
\bottomrule
\end{tabular}
\vspace{4mm}

\begin{itemize}
  \small
  \item Sobre o corpus completo (80 mil notícias): \textbf{0,595}
  \item \textbf{Capta tema, não direção}, e é por isso que os casos aparecem um a um
\end{itemize}
\end{frame}

\begin{frame}{Resultado 3: triagem \textcolor{red!60!black}{\textbf{não}}}
\centering
\begin{tabular}{lc}
\toprule
\textbf{modelo} & \textbf{PR-AUC} \\
\midrule
\textbf{Só volatilidade} & $\mathbf{0{,}542}$ \\
Só contexto & $0{,}538$ \\
Contexto + texto & $0{,}496$ \\
Só texto & $0{,}439$ \\
\bottomrule
\end{tabular}
\vspace{3mm}

\begin{center}
\textbf{Nenhum modelo que lê o texto bate a volatilidade.}
\end{center}
\vspace{2mm}
\footnotesize
Era o contrário do que eu esperava. Verifiquei por três caminhos:
re-teste em condições favoráveis ao texto ($0{,}533$, ainda abaixo),
reamostragem por grupos, e \textbf{nove} definições diferentes do rótulo: a volatilidade ganha
nas nove.
\end{frame}

\begin{frame}{O erro que encontrei em mim próprio}
Eu tinha escrito que a triagem \emph{``quase quadruplica''} a precisão: de $0{,}163$ para $0{,}632$.
\vspace{3mm}

O $0{,}163$ vinha do modelo ``alertar sempre'', que dá a mesma pontuação a tudo.\\
\textbf{Com tudo empatado, a métrica desempata pela ordem do ficheiro}, que está ordenado por data
e por nome da empresa.
\vspace{3mm}
\begin{center}
\large Esse número não media escolher às cegas.\\
\textbf{Media escolher por ordem alfabética.}
\end{center}
\vspace{3mm}
\footnotesize
Ao acaso a sério dá $0{,}379$ $\Rightarrow$ o ganho é de \textbf{1,67$\times$}, não de quase quatro.\\
E uma tabela de \textbf{13 constantes} dá $0{,}662$: bate o modelo treinado.
\end{frame}

% ══ 6. O QUE APRENDI ════════════════════════════════════════════════════════
\begin{frame}{O que aprendi}
\begin{enumerate}
  \item \textbf{A linha de base é tão importante como o modelo.}\\
        \footnotesize O resultado mais útil do trabalho veio de olhar para aquilo contra o qual eu
        estava a comparar.\\[3mm]
  \item \textbf{O sítio onde se põe um modelo importa tanto como o modelo.}\\
        \footnotesize O meu funcionava sozinho e deixou de servir atrás de dois filtros simples, que
        já faziam o trabalho.\\[3mm]
  \item \textbf{A técnica mais simples ganhou três vezes.}\\
        \footnotesize Ao Isolation Forest, ao modelo que lê o texto, e ao modelo implantado.
\end{enumerate}
\vspace{3mm}
\begin{center}
\small A parte difícil não foi escolher a técnica mais avançada.\\
Foi montar a avaliação \textbf{capaz de dizer que ela não era precisa}.
\end{center}
\end{frame}

\begin{frame}{Limitações, ditas antes de me perguntarem}
\begin{itemize}
  \item \textbf{Não fiz estudo com utilizadores.} A fidelidade das explicações está garantida por
        construção; a \emph{utilidade} não foi medida, e não a afirmo.
  \item \textbf{Os critérios de alerta ainda não estão bem calibrados}, e noto-o a usar. Medido:
        havia notícia captada em $88{,}5\%$ dos dias invulgares, e é um \emph{limite superior},
        porque pergunta se existia \emph{uma}, não se chegou \emph{a certa}.
  \item \textbf{Só manchetes}, não o corpo dos artigos.
  \item \textbf{O atraso é da fonte, não da infraestrutura}: $\approx$2,5\,h até detetar,
        1\,s até entregar.
\end{itemize}
\end{frame}

% ══ 7. DEMONSTRAÇÃO ═════════════════════════════════════════════════════════
\begin{frame}{Demonstração}
\begin{center}
\includegraphics[width=0.78\textwidth]{app_v5_overview.png}
\end{center}
\vspace{-2mm}
\begin{center}
\footnotesize O sistema está no ar. \textbf{[gravação da aplicação]}
\end{center}
% ⚠️ AO ALUNO: substituir a imagem acima pela GRAVAÇÃO no dia da defesa.
%    Opção A (mais segura): abrir o vídeo em janela separada e voltar aos slides.
%    Opção B: \usepackage{multimedia} e \movie[...]{...}{demo.mp4} — só funciona em
%             alguns leitores de PDF; testar no portátil da apresentação ANTES.
%    Plano B se não houver rede: esta imagem já mostra o essencial.
\end{frame}

\begin{frame}[plain]
\vspace{1.5cm}
\begin{center}
{\Large \textbf{Obrigado}}\\[6mm]
{\large Perguntas?}\\[10mm]
\footnotesize
Duas respostas de ``sim'', uma de ``não'', e um erro meu encontrado e corrigido.\\
\textbf{É isso que uma avaliação honesta produz.}
\end{center}
\end{frame}

\end{document}
